% H375 complete-chain technical report
% All numerical values in this document are copied from the archived CSV
% summaries under epsilon-active/experiments and are intentionally reported
% with their audit/provenance qualifiers.
\documentclass[UTF8,a4paper,11pt]{ctexrep}

\usepackage[a4paper,margin=2.15cm,headheight=15pt]{geometry}
\usepackage{amsmath,amssymb,bm}
\usepackage{booktabs,longtable,tabularx,array,multirow}
\usepackage{graphicx,float}
\usepackage{xcolor}
\usepackage{tikz}
\usetikzlibrary{arrows.meta,positioning,fit,calc,shapes.geometric}
\usepackage{enumitem}
\usepackage{fancyhdr}
\usepackage{caption}
\usepackage{hyperref}
\usepackage{url}

\hypersetup{
  colorlinks=true,
  linkcolor=blue!55!black,
  citecolor=blue!55!black,
  urlcolor=blue!55!black,
  pdftitle={H375 完整求解链条技术报告},
  pdfauthor={epsilon-active experiment}
}

\graphicspath{{figures/}}
\setlength{\parindent}{2em}
\setlength{\parskip}{0.25em}
\setlength{\emergencystretch}{3em}
\setlist{nosep,leftmargin=2.3em}
\renewcommand{\arraystretch}{1.18}
\captionsetup{font=small,labelfont=bf}

\definecolor{stageblue}{RGB}{0,114,178}
\definecolor{stageorange}{RGB}{230,159,0}
\definecolor{stagegreen}{RGB}{0,158,115}
\definecolor{stagered}{RGB}{213,94,0}
\definecolor{stagepurple}{RGB}{204,121,167}
\definecolor{stagegray}{RGB}{95,95,95}
\definecolor{softblue}{RGB}{232,241,248}
\definecolor{softorange}{RGB}{255,240,216}
\definecolor{softgreen}{RGB}{230,244,238}
\definecolor{softpurple}{RGB}{245,234,242}

\pagestyle{fancy}
\fancyhf{}
\fancyhead[L]{H375 完整求解链条}
\fancyhead[R]{精确非光滑 GP 技术报告}
\fancyfoot[C]{\thepage}

\newcommand{\HPWL}{\operatorname{HPWL}}
\newcommand{\Overflow}{\operatorname{Overflow}}
\newcommand{\clip}{\operatorname{clip}}
\newcommand{\R}{\mathcal{R}}
\newcommand{\M}{\mathcal{M}}
\newcommand{\Fset}{\mathcal{F}}
\newcommand{\Eset}{\mathcal{E}}
\newcommand{\Bset}{\mathcal{B}}
\newcommand{\dd}{\mathrm{d}}

\title{H375 完整求解链条\ 技术报告\\
\large 精确矩形重叠密度、epsilon-active HPWL 与容量恢复实验}
\author{epsilon-active 独立 C++ 全局布局框架}
\date{2026 年 8 月 28 日}

\begin{document}

\maketitle
\thispagestyle{empty}

\begin{abstract}
本文对 adaptec1 数据集上的 H375 完整继承求解链条进行可复现记录。布局变量是可移动标准
单元的连续 $(x,y)$ 坐标，最终质量指标始终由带 pin offset 的精确 max--min HPWL 与精确
矩形--网格交叠面积计算。固定宏器件以不可移动的 bin 占用计入容量账本，\texttt{terminal\_NI}
端口不计入物理容量；整个框架不执行行吸附、合法化、详细布局或合法性检查。

链条从 fresh center--Gaussian raw 初始化开始，经 H219 HPWL-only seed、128 网格 H219
静电同伦、512 网格 H221 细化，再进入 H252--H257 的 exact 非光滑恢复，最后经过 H372
surplus-only hierarchical capacity assignment、十轮 exact recovery 和 H375 五轮等形状交换。
H375 终点为 $85.999318$M HPWL、$6.999990\%$ 精确 overflow。与 DREAMPlace 论文 Table II
的 adaptec1 GP 基线 $73.22$M 相比，HPWL 高 $17.453\%$；按 checkpoint handoff 的完整
链条时间为 $236.40$ s，而论文 40-thread GP 时间为 $67$ s。因此，本文不宣称达到了论文
的 1\% HPWL 和两倍运行时间门槛。需要特别区分的是：H219/H221 的早期正 $\mu$ 阶段确实
使用了 DCT/Poisson 静电辅助；H252 之后的报告尾部才是严格的 exact overlap、exact HPWL
非光滑求解。该区分是本报告的主要审计结论。
\end{abstract}

\noindent\textbf{关键词：} 全局布局；非光滑优化；精确矩形重叠；epsilon-active；固定宏容量；
hierarchical assignment；H375

\tableofcontents
\listoffigures

\chapter{报告范围与结论}

\section{问题背景}
挑战 5 要求在 ISPD2005 混合尺寸网表上完成 global placement（GP）。本报告只讨论连续
GP，不生成 legalized placement。研究目标是将网络线长和密度约束写成可审计的非光滑模型，
并检验一条多阶段求解链是否能同时满足质量和时间门槛。H375 是当前 a1 分支中记录的完整
继承链，而不是从某个历史合法化结果重新开始的单次运行。

\section{锁定约束}
以下约束在代码、实验协议和结果审计中保持一致：
\begin{enumerate}
  \item HPWL 使用精确的 max--min 形式；epsilon-active 只选择活动片段，不修改 HPWL 值。
  \item 密度使用每个矩形与每个 bin 的精确交叠面积；不使用 Gaussian、WA/LSE、Moreau
        envelope、卷积或其它目标平滑。
  \item 固定宏保持坐标不变，并以其真实交叠面积扣减可移动容量。
  \item 不执行 legalization、row/site snapping、detailed placement 或 legality check。
  \item OpenMP 线程请求限制为 $1$--$40$。H219--H257 的历史 checkpoint 使用 16 线程；
        H372/H375 使用 40 线程。
\end{enumerate}

\section{核心结论一览}
\begin{table}[H]
\centering
\caption{H375 终点与论文门槛的审计比较。}
\label{tab:gate}
\small
\begin{tabularx}{\textwidth}{>{\raggedright\arraybackslash}p{2.65cm} >{\raggedright\arraybackslash}p{2.75cm} >{\raggedright\arraybackslash}p{2.75cm} >{\raggedright\arraybackslash}X}
\toprule
指标 & 论文/挑战门槛 & H375 & 判定与解释\\
\midrule
adaptec1 HPWL & $\le 73.9522$M（基线 $+1\%$） & $85.999318$M & 未通过；相对基线劣化 $17.453\%$\\
精确 overflow & $<7\%$ & $6.999990\%$ & 通过（浮点容差内）\\
GP 迭代 & $\le 5000$ & 各 GP 进程最大日志索引 1199/399；尾部用 22 次恢复 sweep & 通过单进程口径；恢复 sweep 不是 Adam iteration\\
运行时间 & $\le 134$ s（$2\times67$ s） & $236.40$ s 完整 handoff 链 & 未通过；仅 H257 之后为 $102.63$ s\\
合法化 & 不纳入本任务 & 未执行 & 符合范围\\
\bottomrule
\end{tabularx}
\end{table}

\section{重要的合规边界}
H375 的最终评价目标是精确非光滑目标，且 H252 之后所有候选均由 exact HPWL 和 exact
overlap 重新审计。然而，H219 和 H221 的正 $\mu$ 同伦阶段调用了 DCT Poisson 静电场，
这与“整条链任何阶段均不得有静电辅助”的最严格解释不一致。因而本文采用以下严谨表述：
H375 是“静电辅助初始化 + exact 非光滑尾部”的完整实验链；它不能被写成从 raw 到终点
全程纯 exact overlap 的证据。若将静电也列为禁用项，则 H219/H221 应被替换，当前 H375
结果不满足该附加约束。

\chapter{数学模型与精确审计}

\section{布局变量和矩形几何}
设放置区域为
\[
  \R=[x_l,x_h]\times[y_l,y_h],
\]
可移动实例集合为 $\M$，固定宏集合为 $\Fset$。实例 $i$ 的连续坐标为 $(x_i,y_i)$，
宽高为 $(w_i,h_i)$，其矩形写为
\[
  R_i(x_i,y_i)=
  [x_i-\tfrac12w_i,x_i+\tfrac12w_i]
  \times
  [y_i-\tfrac12h_i,y_i+\tfrac12h_i].
\]
坐标在 oracle 内被限制在放置区域边界以内，但不被吸附到离散 row 或 site。

\section{精确 HPWL}
对 net $e\in\Eset$，pin $p$ 所属实例为 $i(p)$，Bookshelf pin offset 为
$(\Delta x_p,\Delta y_p)$，则
\[
  X_p=x_{i(p)}+\Delta x_p,\qquad
  Y_p=y_{i(p)}+\Delta y_p.
\]
精确线长为
\begin{equation}
  W(x,y)=\sum_{e\in\Eset}w_e\left[
  \max_{p\in e}X_p-\min_{p\in e}X_p+
  \max_{p\in e}Y_p-\min_{p\in e}Y_p\right].
  \label{eq:hpwl}
\end{equation}
式 \eqref{eq:hpwl} 是所有报告表格中 HPWL 的唯一数值定义。即使计算方向采用活动集合，
前向值仍由 max/min 直接求得。

\section{精确矩形--bin overlap}
将 $\R$ 划分为 $N_x\times N_y$ 个矩形 bin，bin $b$ 的边界为
\([b_x^-,b_x^+]\times[b_y^-,b_y^+]$，面积为 $A_b$。实例 $i$ 对 bin 的一维截长为
\[
 \ell^x_{ib}=\left[\min(x_i+\tfrac12w_i,b_x^+)
 -\max(x_i-\tfrac12w_i,b_x^-)\right]_+,
\]
\[
 \ell^y_{ib}=\left[\min(y_i+\tfrac12h_i,b_y^+)
 -\max(y_i-\tfrac12h_i,b_y^-)\right]_+,
 \qquad a_{ib}=\ell^x_{ib}\ell^y_{ib}.
\]
固定宏占用和总占用分别为
\begin{equation}
 u_b^{\rm fix}=\sum_{j\in\Fset\setminus\{\texttt{terminal\_NI}\}}
 |R_j\cap B_b|,
 \qquad
 u_b=u_b^{\rm fix}+\sum_{i\in\M}a_{ib}.
 \label{eq:occupancy}
\end{equation}
目标密度为 $\rho_t$（实验中为 1），因此固定宏感知的可移动容量是
\begin{equation}
 C_b^{\rm mov}=\left[\rho_t A_b-u_b^{\rm fix}\right]_+.
 \label{eq:movcap}
\end{equation}
实现先把固定矩形写入 immutable occupancy map，再叠加 movable rectangles；这使宏器件
影响直接进入每一个相交 bin，而不是通过额外的软障碍项近似。

\section{精确 overflow 和平方超容量}
报告的 overflow 为
\begin{equation}
 O(x,y)=\frac{\displaystyle\sum_{b\in\Bset}[u_b-\rho_tA_b]_+}
 {A_{\rm mov}},
 \qquad A_{\rm mov}=\sum_{i\in\M}w_ih_i.
 \label{eq:overflow}
\end{equation}
为实现难题 5 中的平方 density excess，GP 阶段还计算
\begin{equation}
 E_{\rm ex}(x,y)=
 \frac12\sum_{b\in\Bset}A_b
 \left[\frac{u_b}{A_b}-\rho_t\right]_+^2
 =\frac12\sum_{b\in\Bset}
 \frac{[u_b-\rho_tA_b]_+^2}{A_b}.
 \label{eq:excess}
\end{equation}
平方项只改变 penalty 的数值权重，不改变 occupancy 和 overflow 的几何定义。对一个处于
固定 active piece 的 bin，精确 piecewise 导数为
\begin{equation}
 \frac{\partial E_{\rm ex}}{\partial x_i}
 =\sum_b e_b\frac{\partial a_{ib}}{\partial x_i},
 \qquad
 \frac{\partial a_{ib}}{\partial x_i}
 =\frac{\partial\ell^x_{ib}}{\partial x_i}\ell^y_{ib},
 \quad
 e_b=\left[\frac{u_b}{A_b}-\rho_t\right]_+.
 \label{eq:densitygrad}
\end{equation}
在矩形完整落于同一 bin 的内部区域，$\partial a_{ib}/\partial x_i$ 和
$\partial a_{ib}/\partial y_i$ 可以同时为零；这正是 extreme subgradient sparsity 的来源。
平方 excess 会放大已经 active 的超载 bin，却不会自动产生穿越 breakpoint 的方向。

\section{epsilon-active HPWL 方向}
当 net 的极值 pin 唯一时，普通 max/min 方向会集中在少数 pin。对 $x$ 方向，以最大端
为例，代码使用
\[
 a^{\max}_{p,x}=\left[1-\frac{x_e^{\max}-X_p}{\epsilon}\right]_+^q,
\]
对最小端定义类似的 $a^{\min}_{p,x}$，再令
\begin{equation}
 g_{p,x}^{(e)}=w_e\left(
 \frac{a^{\max}_{p,x}}{\sum_{r\in e}a^{\max}_{r,x}}
 -\frac{a^{\min}_{p,x}}{\sum_{r\in e}a^{\min}_{r,x}}
 \right).
 \label{eq:epshpwl}
\end{equation}
$y$ 方向完全相同。H219--H257 使用 $\epsilon=125$、$q=4$；精确尾部的前向 HPWL
仍然使用式 \eqref{eq:hpwl}。因此 epsilon 是次梯度选择器，而不是 WA、LSE 或其它平滑
wirelength 模型。

对 overlap 的边界方向，代码可以在 $\epsilon>0$ 时把相邻 affine pieces 纳入候选；
H252--H375 的正式 exact 审计将 \texttt{density\_epsilon}=0，并通过 breakpoint、
coordinate 和 group proposals 生成有限候选。每个候选都重新计算式 \eqref{eq:hpwl} 和
式 \eqref{eq:overflow}，不把伪方向写成真实次梯度。

\chapter{优化器、lambda 与同伦控制}

\section{可扩展优化器接口}
框架将优化器与两个 oracle 解耦。优化器只实现
\[
 \Delta z_k=\operatorname{Solver}(g_k,\eta_k,\text{state}_k),
 \qquad z=(x_1,\ldots,x_{|\M|},y_1,\ldots,y_{|\M|}),
\]
并支持 Adam、AMSGrad、AdaGrad、Heavy-ball 和 SGD。H375 链的 GP 子阶段选用 Adam，
而 exact recovery 使用候选审核，不依赖某个特定的连续优化器。这种边界使后续替换
solver 不必改写 HPWL 或 overlap oracle。

\section{DreamPlace 风格自适应 lambda}
在有 density direction 的 GP 阶段，首先用初始梯度范数校准物理尺度：
\begin{equation}
 \lambda_0=s_\lambda
 \frac{\|g_W\|_1}{\max(\|g_E\|_1,10^{-30})},
 \qquad
 \lambda_k=\lambda_0c_k.
 \label{eq:lambda0}
\end{equation}
其中 $g_E$ 来自式 \eqref{eq:densitygrad} 的 exact piecewise 方向。DreamPlace policy
根据相邻 HPWL 变化更新控制量 $c_k$：若 $\Delta W_k<0$，
\begin{equation}
 c_{k+1}=c_k\,1.05\max(0.9999^k,0.98),
 \label{eq:lambdadec}
\end{equation}
否则使用
\begin{equation}
 c_{k+1}=c_k\,\clip\left(
 1.05^{1-\Delta W_k/\Delta W_{\rm ref}},0.95,1.05\right).
 \label{eq:lambdainc}
\end{equation}
当 overflow 仍明显高于目标时，控制量还受到上限和 deadband 的约束。需要注意，
H219 HPWL-only seed 的 density scale 为零；其中日志里的 lambda control 虽然增加，
实际密度方向仍未进入坐标更新。

\section{\texorpdfstring{正 $\mu$ 静电同伦的定义与审计}{正 mu 静电同伦的定义与审计}}
为解释 H219/H221 的早期长程搬运，定义 DCT Poisson 辅助项
\begin{equation}
 -\nabla^2\phi=\rho,\qquad
 E_{\rm el}=\frac12\int_{\R}\rho\phi\,\dd A,qquad
 \bm{E}=-\nabla\phi.
 \label{eq:electro}
\end{equation}
离散实现采用 DREAMPlace/ePlace 风格的 DCT、IDCT 和 DST/IDST 频域算子。阶段目标可写成
归一化形式
\begin{equation}
 F_{\mu,\lambda}=
 \frac{W}{W_0}
 +\lambda\frac{E_{\rm ex}}{E_0}
 +\mu\left(\frac{E_{\rm el}}{E_{{\rm el},0}}+E_{\rm anchor}\right).
 \label{eq:homotopy}
\end{equation}
在 H219/H221 中 $\mu$ 从正值衰减，$\lambda$ 在 overflow 下降和 HPWL 响应后接管；
报告尾部把物理 $\mu$ 强制为零。式 \eqref{eq:homotopy} 不是 H375 最终的非光滑评价
目标，只是已有 checkpoint 的初始化工具，故必须单独披露。

\chapter{H375 完整链条与 checkpoint provenance}

\section{流程总览}
图 \ref{fig:flow} 用颜色区分初始化、正 $\mu$ 同伦和 exact 尾部。箭头表示“下一阶段
读取上一阶段坐标文件”，不是把各阶段的 objective 直接相加。

\begin{figure}[H]
\centering
\begin{tikzpicture}[
  stage/.style={draw,rounded corners=2pt,align=center,text width=3.25cm,
                minimum height=0.86cm,font=\small,inner sep=3pt},
  arr/.style={-{Latex[length=2mm]},line width=0.65pt},
  node distance=0.65cm and 0.48cm
]
\node[stage,fill=softblue,draw=stageblue] (raw) {fresh raw\\center--Gaussian\\$53.2148$M / $99.966\%$};
\node[stage,fill=softblue,draw=stageblue,right=of raw] (seed) {H219 HPWL seed\\Adam, $\lambda=0$\\$43.0348$M / $96.056\%$};
\node[stage,fill=softpurple,draw=stagepurple,right=of seed] (coarse) {H219 coarse\\128$\times$128 DCT\\handoff $60.6175$M / $71.708\%$};
\node[stage,fill=softpurple,draw=stagepurple,right=of coarse] (fine) {H221 fine\\512$\times$512\\selector $109.426$M / $7.752\%$};
\draw[arr] (raw)--(seed); \draw[arr] (seed)--(coarse); \draw[arr] (coarse)--(fine);

\node[stage,fill=softorange,draw=stageorange,below=1.05cm of raw] (r252) {H252\\15\% exact recovery\\$94.5485$M / $15\%$};
\node[stage,fill=softorange,draw=stageorange,right=of r252] (bridge) {H253--H254\\bridge + retighten\\$89.5493$M / $6.557\%$};
\node[stage,fill=softgreen,draw=stagegreen,right=of bridge] (polish) {H255--H257\\7\% exact polish\\$87.1195$M / $7\%$};
\node[stage,fill=softgreen,draw=stagegreen,right=of polish] (assign) {H372\\surplus-only cuts\\$93.7055$M / $2.306\%$};
\draw[arr] (fine.south) |- (r252.north); \draw[arr] (r252)--(bridge)--(polish)--(assign);

\node[stage,fill=softgreen,draw=stagegreen,below=1.05cm of assign] (rec) {H372 recovery $\times10$\\exact candidates\\$86.1896$M / $7\%$};
\node[stage,fill=softorange,draw=stageorange,left=of rec] (swap) {H375 exchange $\times5$\\equal-shape, net-aware\\$85.9993$M / $7\%$};
\draw[arr] (assign.south)--(rec.north); \draw[arr] (rec)--(swap);
\node[draw=none,below=0.18cm of coarse,font=\footnotesize,text=stagepurple]
  {静电辅助初始化（正 $\mu$）};
\node[draw=none,below=0.18cm of bridge,font=\footnotesize,text=stagegreen]
  {exact 非光滑尾部（$\mu=0$）};
\end{tikzpicture}
\caption{H375 的坐标继承链。每个框中的指标均为该 checkpoint 的 exact audit；H221 的
正 $\mu$ 分支和 H252 之后的 exact 尾部被明确分色。}
\label{fig:flow}
\end{figure}

\section{输入输出关系}
\begin{table}[H]
\centering
\caption{关键 checkpoint 的输入、操作和输出。}
\label{tab:provenance}
\small
\begin{tabularx}{\textwidth}{>{\raggedright\arraybackslash}p{1.45cm} >{\raggedright\arraybackslash}p{3.55cm} >{\raggedright\arraybackslash}X >{\raggedright\arraybackslash}p{2.75cm}}
\toprule
阶段 & 输入坐标 & 核心操作 & 输出坐标/选择规则\\
\midrule
raw & raw benchmark \texttt{adaptec1.pl} & fresh center--Gaussian，seed 219，$\sigma=0.001$ & raw 初始状态\\
H219 seed & raw fresh & HPWL-only Adam，200 updates，density scale 0 & \texttt{global.pl}\\
H219 coarse & H219 seed & 128 网格，正 $\mu$ DCT Poisson 同伦，24$\times$50 steps & 读取 row 1100\\
H221 fine & H219 row 1100 & 512 网格，正 $\mu$ 衰减；$\lambda$ 在 row 60 激活 & row 182 最低 overflow checkpoint\\
H252--H257 & H221 selected & exact breakpoint、axis、compact、net-block；15\% band 后 retighten & H257 \texttt{global.pl}\\
H372 assign & H257 \texttt{last.pl} & fixed-aware position-seeded surplus-only recursive cuts & \texttt{bisect/last.pl}\\
H372 recovery & H372 assignment & exact node/group candidates，7\% cap，10 sweeps &
\texttt{recovery10/}\allowbreak\texttt{last.pl}\\
H375 & H372 recovery10 & 同宽高节点的 net-aware identity exchange，5 sweeps & 最终 \texttt{run/global.pl}\\
\bottomrule
\end{tabularx}
\end{table}

\section{阶段结果总表}
表 \ref{tab:trajectory} 列出绘图脚本从 CSV 读取的完整 checkpoint 轨迹。由于 H252--H375
每一行都是继承坐标后的 exact operator event，不能把它们误读成同一 Adam run 的连续迭代。

\begin{longtable}{p{1.35cm}p{2.45cm}r r p{5.6cm}}
\caption{H375 完整链条的 exact checkpoint 轨迹。HPWL 单位为 M，overflow 为百分数。}\label{tab:trajectory}\\
\toprule
阶段 & checkpoint & HPWL (M) & overflow (\%) & 备注\\
\midrule
\endfirsthead
\multicolumn{5}{c}{\tablename\ \thetable{}（续）}\\
\toprule
阶段 & checkpoint & HPWL (M) & overflow (\%) & 备注\\
\midrule
\endhead
\midrule
\multicolumn{5}{r}{续下页}\\
\endfoot
\bottomrule
\endlastfoot
初始化 & raw fresh & 53.214813 & 99.966213 & center--Gaussian raw benchmark initialization\\
H219 & HPWL seed & 43.034762 & 96.056285 & 200 Adam updates，$\lambda$ 实际无效（scale 0）\\
H219 & coarse row 1100 & 60.617468 & 71.707837 & 128$\times$128 正 $\mu$ 同伦 handoff\\
H221 & fine start & 60.618692 & 77.650357 & 512$\times$512 初始 exact audit\\
H221 & selector row 182 & 109.425511 & 7.751547 & positive-$\mu$ 分支最低 overflow\\
H252 & audit & 109.425511 & 7.803537 & 重新审核的微小 provenance 差异\\
H252 & sweep 1 & 97.813945 & 14.999745 & 15\% relaxed exact recovery\\
H252 & sweep 2 & 94.548507 & 14.999997 & 15\% band endpoint\\
H253 & row 80 & 86.241998 & 13.891233 & moderate exact bridge\\
H254 & selected row 99 & 89.549340 & 6.556721 & strong retightening\\
H255 & sweep 1 & 88.064828 & 6.999984 & 7\% exact candidate polish\\
H255 & sweep 2 & 87.567244 & 6.999999 & 7\% exact candidate polish\\
H255 & sweep 3 & 87.317935 & 6.999977 & 7\% exact candidate polish\\
H255 & sweep 4 & 87.178874 & 7.000000 & 7\% exact candidate polish\\
H257 & sweep 5 & 87.119465 & 6.999958 & fifth exact node sweep\\
H372 & assignment & 93.705492 & 2.306362 & surplus-only hierarchical capacity assignment\\
H372 & recovery 1 & 89.143747 & 6.999985 & exact slack spending\\
H372 & recovery 2 & 87.641415 & 6.999963 & exact slack spending\\
H372 & recovery 3 & 86.972538 & 6.999991 & exact slack spending\\
H372 & recovery 4 & 86.632343 & 6.999971 & exact slack spending\\
H372 & recovery 5 & 86.450273 & 6.999985 & exact slack spending\\
H372 & recovery 6 & 86.343968 & 6.999999 & exact slack spending\\
H372 & recovery 7 & 86.278326 & 7.000000 & exact slack spending\\
H372 & recovery 8 & 86.234833 & 7.000000 & exact slack spending\\
H372 & recovery 9 & 86.207181 & 6.999993 & exact slack spending\\
H372 & recovery 10 & 86.189626 & 6.999990 & exact slack spending\\
H375 & exchange 1 & 86.064268 & 6.999990 & equal-shape exchange\\
H375 & exchange 2 & 86.025509 & 6.999990 & equal-shape exchange\\
H375 & exchange 3 & 86.009906 & 6.999990 & equal-shape exchange\\
H375 & exchange 4 & 86.002560 & 6.999990 & equal-shape exchange\\
H375 & exchange 5 & \textbf{85.999318} & \textbf{6.999990} & final exact audit\\
\end{longtable}

\chapter{阶段一：fresh 初始化与 H219 HPWL seed}

\section{初始化动机}
DREAMPlace 论文观察到中心加小 Gaussian 噪声能够避免昂贵的 bound-to-bound 初始摆放。
本链条遵循相同的初始化思想，但每次实验都从 raw \texttt{.pl} 重新生成坐标，不读取
历史 \texttt{.eplace*.pl} 或 \texttt{.lg.pl}。对 H375，seed 为 219，噪声标准差与版图
宽高的比例为 $10^{-3}$。因此初始单元集中在中心，exact movable overlap 很高：
\[
  W_0=53.214813\text{M},\qquad O_0=99.966213\%.
\]

\section{HPWL-only seed}
第一阶段只令 $\lambda E_{\rm ex}=0$，采用 Adam 和式 \eqref{eq:epshpwl} 的 HPWL
epsilon-active 方向。密度仍被 exact 计算和记录，但不进入坐标更新。200 次更新后
\[
  53.214813\text{M}\longrightarrow43.034762\text{M},
  \qquad
  99.966213\%\longrightarrow96.056285\%.
\]
这一步的作用是让同一 net 的 pin 先形成较紧的拓扑，而不是建立容量可行性。结果也说明
“只优化 HPWL”不能解决中心拥挤；如果此时强行使用逐节点 exact density gradient，很多单元
仍处于同一 bin 内部，方向会为零。

\section{收敛行为}
图 \ref{fig:full} 左侧的初始化/HPWL 区间显示 HPWL 下降与 overflow 几乎脱钩。
日志中的 \texttt{lambda\_control} 在该阶段增加，是控制器的记录值；由于 density weight
scale
为 0，物理坐标更新没有密度力。这一区别是解释 H219 seed 的必要条件。

\chapter{\texorpdfstring{阶段二：H219 粗网格正 $\mu$ 同伦}{阶段二：H219 粗网格正 mu 同伦}}

\section{设置}
H219 读取 H219 seed，在 128$\times$128 网格上运行 24 个 50-step stage。初始控制参数
为 $\mu_{\rm control}=5.36$，物理静电尺度约为 $0.443690$；exact overlap scale
约为 $0.00674558$。$\mu$ 按 $0.99$ 衰减，并在稳定窗口中交给 adaptive controller；
lambda step 为 15、最大 control 为 150，anchor weight 为 $10^{-4}$。每个 iteration
仍计算 exact rectangle overlap 作为日志指标，但早期主要坐标方向来自式 \eqref{eq:electro}。

\section{结果和物理解释}
粗网格初始状态的 exact audit 为
\[
  W_{\mathrm{coarse},0}=43.034762\text{M},\qquad
  O_{\mathrm{coarse},0}=94.7725\%.
\]
row 0 的细节见 \texttt{homotopy\_metrics.csv}。静电场提供跨 bin 的长程排斥，使 overflow
逐步降到 row 1100 的 $71.707837\%$；与此同时 HPWL 上升至 $60.617468$M。
该变化不是“平方 overlap 梯度单独成功”，而是正 $\mu$ 场把大量中心单元整体搬向版图
外侧。

粗网格的 handoff 选在 row 1100，是为了把已产生的全局 spread 交给细网格，而不是把
row 1199 的末端状态直接当作下一阶段。图 \ref{fig:homo} 左上给出完整粗网格轨迹，竖线
标注了该 handoff。

\chapter{阶段三：H221 细网格同伦与 selector 语义}

\section{512 网格切换}
H221 读取 H219 row 1100 坐标并重新建立 512$\times$512 fixed-macro occupancy。因为
bin 尺寸骤减，exact rectangle/bin geometry 被重新计算，row 0 的 overflow 回升到
$77.650357\%$，HPWL 为 $60.618692$M。这一跳是分辨率改变造成的 exact metric 变化，
不是坐标合法化。

参数为 $\mu_{\rm control}=964$、$\mu$ decay $=0.99$、electrostatic gradient scale
$=0.0979419$、overlap scale $=0.000925412$、lambda step $=520$、lambda maximum
$=1950$、initial anchor weight $=10^{-5}$。H221 的 $\lambda$ 在 iteration 0--59
保持为 0；iteration 60 开始由 adaptive controller 激活。图 \ref{fig:homo} 右下显示
在 overlap term 被禁用的早期，overflow 仍随 electrostatic term 明显下降。

\section{row 182 checkpoint}
正 $\mu$ 分支在 row 182 达到日志中的最低 exact overflow：
\[
  W_{182}=109.425511\text{M},\qquad O_{182}=7.751547\%.
\]
它尚未达到 7\%，但已进入后续 exact recovery 可处理的狭窄容量带。继续运行时，HPWL
下降而 overflow 回升；row 349 已为 $96.3050$M / $13.1115\%$ 左右。因此 selector
选择 row 182，而不是把末端 HPWL 较小的状态当作密度可行点。

\section{\texorpdfstring{restore 与 $\mu=0$ 的事实}{restore 与 mu=0 的事实}}
在 row 350，程序恢复 row 182 的坐标，然后把物理 $\mu$ 强制为 0。row 351 的首次
纯 $\mu=0$ 更新立即变成 $107.8691$M / $16.8234\%$；row 399 为 $92.8575$M /
$23.7449\%$。所以 H221 的 $7.751547\%$ 不是一条从 row 349 继续单调下降的纯
non-smooth trajectory，而是“positive-$\mu$ 最低 overflow checkpoint + restore”。
H252 重新读取该文件时得到 $7.803537\%$，存在约 $0.052$ 个百分点的重算差异；本文
保留两套数字并在表 \ref{tab:trajectory} 中标明 provenance。

\chapter{阶段四：H252--H257 的 exact 非光滑尾部}

\section{H252：15\% relaxed exact recovery}
H252 从 H221 selector 开始，先把 overflow cap 放宽到 15\%，以便恢复因 spread 产生的
HPWL 拓扑。候选包括 axis-separated displacement、breakpoint oracle、compact direction、
node move 和最多 16 节点的 net-connected density direction。每个 node 或 group move
都在临时 occupancy 上计算
\[
 \Delta W=W(z+\Delta z)-W(z),\qquad
 \Delta O=O(z+\Delta z)-O(z),
\]
并在 exact cap 下接受。两轮 sweep 的结果为
\[
 109.425511\text{M}/7.803537\%
 \rightarrow97.813945\text{M}/14.999745\%
 \rightarrow94.548507\text{M}/14.999997\%.
\]
此阶段的核心不是修改 density 函数，而是允许暂时使用 15\% slack，让 net-relative
geometry 有机会恢复。

\section{H253：moderate exact bridge}
H253 在 512 网格上使用 Adam，density weight scale $0.5$、net-batch weight 1、
step fraction $0.002$。前向 density 仍是式 \eqref{eq:excess}，\texttt{density\_epsilon}=0；
net-batch 仅作为共享方向的附加搜索项，不改变目标定义。row 80 的 exact endpoint 为
\[
  W=86.241998\text{M},\qquad O=13.891233\%.
\]
这是一个短桥接阶段：它把 relaxed band 中的 HPWL 继续压低，但不强行把 overflow
一次推到 7\%。

\section{H254：strong retightening}
H254 将 density weight scale 提高到 4.75，step fraction 降到 $0.001$，保持 net-batch
weight 1。按照 exact checkpoint 选择规则，row 99 为较好的 retightening 状态：
\[
  W=89.549340\text{M},\qquad O=6.556721\%.
\]
HPWL 相对 H253 暂时上升，是用容量压力换取进入 7\% band；该阶段仍没有静电项、WA/LSE
或平滑 density field。

\section{H255/H257：7\% cap 下的 node polish}
H255 固定 $O\le0.07$，进行四轮 exact node-wise recovery；H257 再做第五轮。候选来自
epsilon-active HPWL、axis steps、exact bin-boundary breakpoints、net centroid 和
support centroid compact directions。接受条件为
\begin{equation}
  \Delta W<0,
  \qquad O(z+\Delta z)\le0.07+\tau,
  \label{eq:feasible_accept}
\end{equation}
并在当前超过 cap 时要求 exact overflow 不增加。结果为
\[
\begin{aligned}
  89.549340\text{M}/6.556721\%
  &\rightarrow88.064828\text{M}
   \rightarrow87.567244\text{M}\\
  &\rightarrow87.317935\text{M}
   \rightarrow87.178874\text{M}\\
  &\rightarrow\boxed{87.119465\text{M}/6.999958\%}.
\end{aligned}
\]
H257 是后续 H372 的输入。这里的“恢复”是 exact candidate search，而不是把零梯度替换
成固定非零导数；breakpoint proposal 只是优化器层的搜索方向，所有提交仍由式
\eqref{eq:hpwl} 和式 \eqref{eq:overflow} 审核。

\chapter{阶段五：H372 hierarchical capacity assignment}

\section{为什么需要 assignment}
H257 已在 7\% band，但很多节点仍以局部碎片方式占据过载 bin。继续逐节点移动时，
exact overlap 的 active set 过于稀疏；如果允许长距离任意跳转，HPWL 会发生不可逆的
topology debt。H372 因此把容量分配提升到版图区域层级，并只移动每个 cut 两侧真正超出
容量的边界 surplus。

\section{位置 seeded 的递归 cut}
在一个递归区域 $Q$ 内沿较长轴切成 $Q_L,Q_R$。固定宏感知容量记为 $C_L,C_R$，当前
可移动面积为 $A_L,A_R$。位置 seeded 意味着节点先按其当前坐标归属一侧，不做全局 axis-rank
或 Hilbert 重排。若 $A_L>C_L$，只从靠近 cut 的边界节点中取最小前缀 $T$，使
\begin{equation}
 A_L-\sum_{i\in T}a_i\le C_L,
 \qquad
 A_R+\sum_{i\in T}a_i\le C_R.
 \label{eq:surplus}
\end{equation}
同理处理右侧。节点在叶区内使用 nearest-capacity 和 HPWL-guided 选择；固定宏不会被
移动，也不被当作普通 movable node。该策略的要点是“保留原有 side topology，只搬 cut
边界的必要面积”，而不是把所有节点均匀铺满整个版图。

\section{H372 assignment 结果}
配置为 512$\times$512、40 threads、leaf bins 8、degree limit 32、FM pass 1、无
axis pre-map。共产生 5833 splits、5834 leaves、最大递归深度 65。一次 assignment
把
\[
 87.119465\text{M}/6.999958\%
  \longrightarrow93.705492\text{M}/2.306362\%.
\]
HPWL 的暂时上升是有意的：assignment 先建立显式容量 slack，再将该 slack 交给 exact
recovery 使用。图 \ref{fig:tail} 的 assignment 标注清楚显示了这一“先降 overflow、
后花 slack 恢复 HPWL”的事件。

\chapter{阶段六：H372 exact recovery 与 H375 exchange}

\section{H372 recovery 的候选机制}
从 assignment checkpoint 出发，H372 recovery10 做十轮 exact recovery。每轮依次尝试：
\begin{enumerate}
  \item node-wise epsilon-active HPWL 方向和轴分离步长；
  \item exact bin-boundary breakpoint；
  \item compact/support-centroid 方向；
  \item degree $\le16$、最多 8 节点的 net-connected rigid/shared density block；
  \item 对每个临时 node/group move 进行 exact HPWL、exact overflow 和 fixed-capacity 审核。
\end{enumerate}
若当前状态已在 7\% cap 内，则候选不能使式 \eqref{eq:overflow} 超过 cap；若当前状态
超过 cap，则候选还需降低 exact overflow。这样做允许非局部但仍可审计的 net-aware group
方向，避免把 density pseudo-gradient 直接冒充模型梯度。

十轮轨迹为
\[
\begin{aligned}
  93.705492&\rightarrow89.143747\rightarrow87.641415\rightarrow86.972538\\
  &\rightarrow86.632343\rightarrow86.450273\rightarrow86.343968\\
  &\rightarrow86.278326\rightarrow86.234833\rightarrow86.207181\\
  &\rightarrow86.189626\ \text{M},
\end{aligned}
\]
末端 overflow 始终在 $6.999990\%$ 附近。前几轮收益较大，后几轮进入 diminishing
returns，说明 7\% band 内的主要损失已经是早期容量 assignment 造成的拓扑代价。

\section{H375 等形状交换}
H375 只在相同 $(w,h)$ 的节点之间交换坐标中心。若节点 $i,j$ 形状相同，交换前后对任意
bin $b$ 的总面积为
\[
 a_{ib}(p_i)+a_{jb}(p_j)
 =a_{jb}(p_i)+a_{ib}(p_j),
 \]
所以
\begin{equation}
 u_b^{\rm after}=u_b^{\rm before},
 \qquad O^{\rm after}=O^{\rm before},
 \label{eq:swapinvariant}
\end{equation}
在浮点容差内严格成立。交换是否接受只由涉及节点 incident nets 的精确 HPWL delta
决定；net-aware 候选半径为 64，候选数 32，exact shortlist 16，degree limit 1000。
因此 H375 是 density-neutral 的拓扑 polish，不是新的 density model。

五轮分别接受 8064、3415、1651、812、429 次交换，总计 14,371 次，得到
\[
 86.189626\text{M}\rightarrow86.064268\rightarrow86.025509
 \rightarrow86.009906\rightarrow86.002560
 \rightarrow\boxed{85.999318\text{M}},
\]
overflow 从 $6.99999030053\%$ 到 $6.99999030053\%$ 不变。单独看，交换只贡献约
$0.190308$M（$0.221\%$）的 HPWL 改善；它不能弥补从 H221 到 H375 的主要 topology debt。

\chapter{收敛曲线与轨迹分析}

\section{完整 checkpoint 曲线}
图 \ref{fig:full} 同时绘出 exact HPWL 和 exact overflow。背景色对应流程阶段，水平虚线
分别是 DREAMPlace a1 HPWL 基线 $73.22$M 和 7\% overflow 目标。曲线揭示三个事实：
\begin{enumerate}
  \item H219 seed 先改善 HPWL，但中心拥挤使 overflow 仍接近 100\%。
  \item H219/H221 正 $\mu$ 同伦承担了主要的全局 spread；fine selector 以较高 HPWL
        换取约 7.75\% overflow。
  \item H372 assignment 形成明显的“HPWL 上跳、overflow 下跳”事件，随后 H372 recovery
        和 H375 exchange 只在窄 band 内逐步收回 HPWL。
\end{enumerate}

\begin{figure}[H]
\centering
\includegraphics[width=\textwidth]{fig_full_chain_convergence.pdf}
\caption{H375 完整继承 checkpoint 曲线。上图为精确 HPWL，下图为精确 overflow；阶段
背景、基线、7\% 目标及关键 checkpoint 均已标注。}
\label{fig:full}
\end{figure}

\section{同伦细节和力项归因}
图 \ref{fig:homo} 给出 H219 粗网格、H221 细网格 selector、$\mu/\lambda$ 控制以及
归一化梯度项。H221 的竖线分别位于 row 60（lambda 激活）、row 182（最低 overflow）和
row 350（恢复并强制 $\mu=0$）。图中 row 351--399 的紫色虚线是已记录的纯 $\mu=0$
trial，而不是被选作 H252 输入的状态。

\begin{figure}[H]
\centering
\includegraphics[width=\textwidth]{fig_homotopy_detail.pdf}
\caption{H219/H221 同伦的详细日志轨迹。右上图区分正 $\mu$ branch、被拒绝的后续 branch
和 restore；右下图显示早期 overlap term 为零时，overflow 仍由 electrostatic term
下降。}
\label{fig:homo}
\end{figure}

\section{exact 尾部曲线}
图 \ref{fig:tail} 放大 H252--H375。H252 的 15\% topology-recovery band 允许 HPWL
从 109.43M 降到 94.55M；H253/H254 交替进行 topology bridge 和容量 retightening；
H372 surplus-only assignment 再次牺牲 HPWL 以创造 capacity slack；H372 recovery 和
H375 exchange 最终回到 7\% band。该图支持“容量分配是主导机制，等形状交换是小幅 polish”
的因果解释。

\begin{figure}[H]
\centering
\includegraphics[width=\textwidth]{fig_exact_tail.pdf}
\caption{H252--H375 exact 尾部轨迹。上图为 HPWL，下图为 overflow；15\% relaxed band、
7\% target、surplus-only assignment 和每个关键 endpoint 均已标注。}
\label{fig:tail}
\end{figure}

\chapter{运行时间、迭代和复现审计}

\section{时间口径}
论文 Table II 在 40-core Intel E5-2698 v4 上报告 adaptec1 GP 为 67 s，因此挑战门槛
取 134 s。本机早期 H219--H257 checkpoint 运行使用 16 threads，H372/H375 使用 40
threads；不同线程数和不同硬件意味着这里不能声称严格的同机 speedup。为避免把增量时间
误当端到端时间，本报告按 checkpoint handoff 累加：
\begin{equation}
\begin{aligned}
  T_{\rm full}
  &=5.800+64.440+19.889+8.740+4.911+7.614\\
  &\quad+19.090+3.288+2.059+62.942+37.625\\
  &=236.398\text{ s}.
\end{aligned}
 \label{eq:runtime}
\end{equation}
其中 H221 使用 row-350 handoff 的 19.889 s；完整 H221 selector 进程 wall time 为
22.843 s。若把完整进程和所有文件 I/O 均纳入，端到端估计约为 239.35 s，结论不会改变。

\begin{figure}[H]
\centering
\includegraphics[width=\textwidth]{fig_runtime_breakdown.pdf}
\caption{H375 阶段时间 waterfall。虚线为论文 67 s 与两倍门槛 134 s；图中同时解释了
为什么历史记录中的约 103 s 只能代表 H257 之后的增量链，不能代表从 raw 初始化开始的
完整运行。}
\label{fig:runtime}
\end{figure}

\begin{table}[H]
\centering
\caption{H375 各阶段的 checkpoint-handoff 时间账本。}
\label{tab:runtime}
\small
\begin{tabular}{p{5.1cm}r r p{4.0cm}}
\toprule
阶段 & 时间 (s) & 累计 (s) & 口径\\
\midrule
H219 HPWL seed & 5.800 & 5.800 & 16 threads，Adam\\
H219 coarse 至 row 1100 & 64.440 & 70.240 & 16 threads，正 $\mu$ DCT\\
H221 fine 至 restore & 19.889 & 90.129 & row-350 handoff；完整进程 22.843 s\\
H252 cap 15\% & 8.740 & 98.869 & exact recovery\\
H253 bridge & 4.911 & 103.780 & Adam bridge\\
H254 retighten & 7.614 & 111.394 & Adam retighten\\
H255 polish & 19.090 & 130.483 & 四轮 exact recovery\\
H257 polish & 3.288 & 133.772 & 第五轮 exact recovery\\
H372 assignment & 2.059 & 135.831 & 40 threads，surplus-only cuts\\
H372 recovery $\times10$ & 62.942 & 198.773 & 40 threads，exact candidates\\
H375 exchange $\times5$ & 37.625 & \textbf{236.398} & 40 threads，equal-shape\\
\bottomrule
\end{tabular}
\end{table}

\section{迭代与 sweep 口径}
H219 seed 使用 200 个 Adam update；H219 coarse 的 handoff index 为 1100（完整日志到
1199）；H221 selector process 到 row 350（完整日志到 399）。H253 和 H254 分别记录到
row 80 和 selected row 99。H252、H255、H257、H372 recovery、H375 exchange 则是
exact operator sweeps，每轮包含大量临时候选和 exact audits，不应伪装成同一种“梯度迭代”。
如果只累加 nominal GP update indices，到 H254 约为
\[
 200+1100+350+80+99=1829<5000,
\]
而 exact 尾部另有 $2+4+1+10+5=22$ 个 sweep。该双重账本比把每个 candidate move 都算成
一次 optimizer iteration 更可复现。

\section{数据和复现命令}
图表由
\url{D:/codex_project/HUAWEI_EDA/epsilon-active/report/h375/figures/gen_h375_figures.py}
生成，输入文件包括：
\begin{itemize}
  \item \url{experiments/h219_a1_hpwl_seed/global_metrics.csv}；
  \item \url{experiments/h219_a1_coarse/homotopy_metrics.csv}；
  \item \url{experiments/h221_a1_fine/homotopy_metrics.csv}；
  \item \url{experiments/h252_a1_relaxed2_net_recovery/recovery_metrics.csv}；
  \item \url{experiments/h253_a1_relaxed2_bridge80/global_metrics.csv}；
  \item \url{experiments/h254_a1_relaxed2_strong475/summary.txt}；
  \item \url{experiments/h255_a1_relaxed2_four_node_polish/recovery_metrics.csv}；
  \item \url{experiments/h257_a1_fifth_node_line4/recovery_metrics.csv}；
  \item \url{experiments/h372_a1_surplus_only_no_axis_map/bisect/bisection_metrics.csv}；
  \item \url{experiments/h372_a1_surplus_only_no_axis_map/recovery10/recovery_metrics.csv}；
  \item \url{experiments/h375_a1_surplus_recovery10_exchange/run/swap_recovery_metrics.csv}。
\end{itemize}
在报告目录下重新生成图表的命令为：
\begin{verbatim}
python figures/gen_h375_figures.py
\end{verbatim}
使用指定 TeX Live 2023 编译：
\begin{verbatim}
D:\latex\101\texlive\2023\bin\windows\xelatex.exe ^
  -interaction=nonstopmode -halt-on-error ^
  -output-directory D:\codex_project\HUAWEI_EDA\epsilon-active\report\h375 ^
  D:\codex_project\HUAWEI_EDA\epsilon-active\report\h375\h375_technical_report.tex
\end{verbatim}

\chapter{讨论：哪些机制真正有效}

\section{容量下降的主导机制}
从轨迹和消融证据看，H375 的 overflow 下降不是由单一的 Adam density gradient 完成的。
主要贡献按作用层次可归纳为：
\begin{enumerate}
  \item H219/H221 的正 $\mu$ electrostatic continuation 提供跨大量 fine bins 的长程
        spread，克服了 raw center seed 的零梯度盆地；
  \item H252--H254 的 15\% relaxed band 和 exact squared-excess bridge 恢复部分 net
        topology，同时将密度逐渐接管；
  \item H372 surplus-only hierarchical cuts 首次显式利用 fixed-macro-aware residual
        capacity，只搬 cut-boundary surplus，产生了 $2.306\%$ overflow checkpoint；
  \item H372 exact recovery 将 assignment 产生的容量余量花回 HPWL，回到 7\% band；
  \item H375 equal-shape exchange 在不改变任何 bin occupancy 的条件下做最后的小幅线长
        改善。
\end{enumerate}
因此，“最关键的容量机制”是 topology-preserving hierarchical capacity assignment，
而“最关键的 raw basin 激活机制”是早期正 $\mu$ 长程 continuation。等形状交换是有效但
非主导的 final polish。

\section{为什么 HPWL 仍高于 baseline}
H375 的 HPWL gap 主要在 fine-grid 可行化之前形成，而不是在最后 7\% band 内由零梯度
造成。H221 把 $60.62$M 的粗网格状态扩散成约 $109.43$M 的 fine-grid checkpoint；
H252--H257 虽恢复到 $87.12$M，H372 assignment 又暂时升至 $93.71$M。后续 exact
recovery 和 exchange 只能在不超过 7\% 的窄可行带内移动，不能重新构造已经被 cut 分配
改变的全局 net topology。单纯增大 lambda、延长 node sweep 或扩大交换半径，会进一步
消耗 HPWL，而不会消除这个早期 topology debt。

\section{零梯度问题的边界}
若节点矩形严格位于单一 bin 内，exact overlap 在该局部区域是常数，真实局部次梯度确实
可以为零。以下做法被允许作为优化器层的候选生成器，但不被称为真实次梯度：
\begin{itemize}
  \item exact breakpoint oracle：计算到相邻 bin 边界的有限位移；
  \item net-aware shared direction：对同一 net 的节点共享刚性或受限弹性方向；
  \item exact finite directional audit：以有限候选步长直接比较 $\Delta O$ 和 $\Delta W$；
  \item controlled noise：仅在 zero-gradient 节点上提出可审计的边界候选。
\end{itemize}
这些方向不修改式 \eqref{eq:occupancy}--\eqref{eq:overflow}，最终接受仍由 exact objective
审核。H405--H410 的单机制实验也显示，breakpoint、net-share 和 finite oracle 能改善
局部搜索，却不能独立把 raw overflow 从约 100\% 送到可行域；因此它们应作为 assignment
之后的 recovery operator，而不是容量桥本身。

\section{网格分辨率和固定宏}
细网格提高了 occupancy 解析度，但同时增加了 active-set 断点数量和零梯度区域。简单地
把 128 网格坐标一次性切换到 512 网格，会产生 7.75\% 的 fine-grid overflow 跳变；H221
的分阶段 handoff 只能缓解，不能消除这一分辨率差异。未来若要进一步改善，应该在每个
分辨率上执行 exact capacity assignment，并把 fixed macro occupancy 和 active bins 显式
prolongate，而不是只复制 electrostatic field。

\chapter{局限性与后续工作}

\section{当前结果的限制}
\begin{enumerate}
  \item H375 只在 adaptec1 上完成；a2--a4 尚未获得相同质量和 runtime 证据，不能外推。
  \item 完整链条包含正 $\mu$ 静电辅助，若挑战解释为全阶段禁止静电，则需重新设计 raw-to-
        capacity continuation。
  \item H375 的 $17.453\%$ HPWL 劣化超过 $1\%$ 门槛；$236.40$ s 也超过 $134$ s。
  \item 所有 overflow 数字是 GP 连续矩形的密度审计，不代表无重叠的合法化版图。
  \item 16-thread 与 40-thread 阶段混用，运行时间只作本实验的 checkpoint 账本，不作跨硬件
        的性能声明。
\end{enumerate}

\section{最有价值的下一步}
在保持最终 exact non-smooth model 的前提下，优先级应为：
\begin{enumerate}
  \item 设计不依赖静电的多分辨率 exact capacity continuation：粗层只生成容量目标，细层
        采用 fixed-macro-aware residual capacity 和 topology-preserving cut；
  \item 在 cut assignment 中加入全局 exact HPWL trust-region 和 net-connected commodity
        的联合 destination assignment，避免一次性产生不可逆 topology debt；
  \item 只在 assignment 后启用 breakpoint/net-share/finite-oracle recovery，并对每个
        candidate 做 exact line search；
  \item 用同一 40-thread 口径重新测量，减少早期 16-thread handoff 的时间不确定性。
\end{enumerate}
这些方向改进的是搜索路径，而不是把 overlap 函数平滑化；最终报告指标仍应由式
\eqref{eq:hpwl} 和式 \eqref{eq:overflow} 计算。

\chapter{结论}

本文已经把 H375 的完整坐标继承链、数学模型、控制策略、exact 尾部算子、收敛曲线和运行
时间逐项记录。实验表明，H375 能在不合法化的情况下把 adaptec1 的精确 overflow 压到
$6.999990\%$，并通过 H372 surplus-only assignment 与 H375 density-neutral exchange
保持固定宏容量和 net-aware HPWL 审计。然而，最终 HPWL 为 $85.999318$M，比 DREAMPlace
GP 基线高 $17.453\%$，完整链条约 $236.40$ s，也未通过两倍时间门槛。

最重要的技术结论有两点。第一，raw center seed 的 extreme subgradient sparsity 不能
靠普通逐节点 Adam 自行解决；需要长程 continuation 或显式 hierarchical capacity bridge
激活新的 exact pieces。第二，breakpoint、net-aware batch 和 equal-shape exchange 可以
作为优化器层的搜索方向，但它们不改变最终非光滑模型；真正决定 overflow 可行化的是
fixed-macro-aware、surplus-only 的区域容量分配。后续工作必须在这一桥接阶段降低 topology
debt，才能接近论文的 HPWL 质量和时间目标。

\appendix
\chapter{参数与审计文件索引}

\begin{table}[H]
\centering
\caption{H375 链条的关键锁定参数。}
\label{tab:params}
\small
\begin{tabular}{p{4.3cm}p{10.2cm}}
\toprule
参数组 & 数值/说明\\
\midrule
raw initialization & fresh center--Gaussian；seed 219；sigma ratio 0.001\\
HPWL oracle & exact pin-offset max--min；$\epsilon=125$；active power 4\\
density oracle & exact rectangle/bin overlap；512$\times$512 fine grid；target density 1\\
fixed macro & immutable occupancy；除 \texttt{terminal\_NI} 外均扣减容量\\
H219 coarse & 128$\times$128；24 stages$\times$50 steps；$\mu_{\rm control}=5.36$；decay 0.99\\
H221 fine & 512$\times$512；8 stages$\times$50 steps；$\mu_{\rm control}=964$；lambda max 1950\\
H252 & 15\% cap；2 sweeps；axis/breakpoint/compact/net-block\\
H253 & density scale 0.5；net-batch weight 1；step 0.002\\
H254 & density scale 4.75；net-batch weight 1；step 0.001；selected row 99\\
H255/H257 & 4+1 exact sweeps；7\% cap；line-search candidates exact audited\\
H372 assignment & surplus-only；leaf bins 8；no axis pre-map；40 threads\\
H372 recovery & 10 sweeps；degree $\le16$；max 8-node blocks；7\% cap\\
H375 exchange & 5 sweeps；radius 64；32 candidates；shortlist 16；equal shapes only\\
\bottomrule
\end{tabular}
\end{table}

\chapter{参考文献}
\begin{thebibliography}{9}
\bibitem{dreamplace}
Y. Lin, Z. Jiang, J. Gu, W. Li, S. Dhar, H. Ren, B. Khailany, and D. Z. Pan,
``DREAMPlace: Deep Learning Toolkit-Enabled GPU Acceleration for Modern VLSI Placement,''
\emph{IEEE Transactions on Computer-Aided Design of Integrated Circuits and Systems},
vol. 40, no. 4, pp. 748--761, 2021. DOI: 10.1109/TCAD.2020.3003843.

\bibitem{eplace}
J. Lu, P. Chen, C.-C. Chang, L. Sha, D. Z. Pan, and C. Y. Hsu,
``ePlace: Electrostatics-Based Placement Using Fast Fourier Transform and Nesterov's Method,''
\emph{ACM Transactions on Design Automation of Electronic Systems}, vol. 20, no. 2, 2015.

\bibitem{ispd}
C. J. Alpert, ``The ISPD2005 placement contest and benchmark suite,''
in \emph{Proceedings of the ACM International Symposium on Physical Design}, 2005.

\bibitem{challenge}
华为 EDA 专题，\emph{第148期 -- EDA 专题第四期}，难题 5，项目约束与数据说明，2026。

\bibitem{implementation}
epsilon-active 独立实现，\url{D:/codex_project/HUAWEI_EDA/epsilon-active}，
源代码模块：
\texttt{hpwl.cpp}、\texttt{density.cpp}、\texttt{lambda\_controller.cpp}、
\texttt{bisection.cpp}、\texttt{recovery.cpp} 和 \texttt{swap\_recovery.cpp}。
\end{thebibliography}

\end{document}
